\documentclass[11pt]{amsart}
\usepackage{amscd,amsmath,amssymb,amsfonts,verbatim}
\usepackage[cmtip, all]{xy}
%\usepackage[letterpaper, left=3cm, right=3cm, top = 2.5cm, bottom = 2.5cm ]{geometry}
\usepackage{comment}
\usepackage{tikz-cd}
\usepackage{hyperref} 




\setlength{\textwidth}{5.8in}
%\setlength{\textheight}{9.2in} % Space saving measures
%\setlength{\textheight}{8.5in}
\setlength{\textheight}{8.7in}
\setlength{\topmargin}{-0.0in}
\renewcommand{\baselinestretch}{1}
\setlength{\oddsidemargin}{.25in}
\setlength{\evensidemargin}{.25in}

%Remove in the final version
%\usepackage[color, notref, notcite]{showkeys}     % refs and labels
%\definecolor{refkey}{gray}{.5}   % graylevel for refs
%\definecolor{labelkey}{gray}{.5} % graylevel for labels
%\definecolor{Red}{rgb}{1,0,0}
\renewcommand\qedsymbol{$\blacksquare$}


\newtheorem{thm}{Theorem}[section]
\newtheorem{prop}[thm]{Proposition}
\newtheorem{lem}[thm]{Lemma}
\newtheorem{cor}[thm]{Corollary}
\newtheorem{conj}[thm]{Conjecture}
\theoremstyle{definition}
\newtheorem{hyp}{Hypothesis}[]


\newtheorem{defn}[thm]{Definition}
%\theoremstyle{remboldstyle}
\newtheorem{rembold}[thm]{Remark}
%\theoremstyle{remark}
%\newtheorem{remk}[thm]{Remark}
\newtheorem{remks}[thm]{Remarks}
\newtheorem{remk}[thm]{Remark}
\newtheorem{exercise}[thm]{Exercise}
\newtheorem{exm}[thm]{Example}
\newtheorem{exms}[thm]{Examples}
\newtheorem{notat}[thm]{Notation}
\numberwithin{equation}{section}
\numberwithin{equation}{section}
\newtheorem{ques}[thm]{Question}






\newenvironment{Def}{\begin{defn}}%
{\hfill$\square$\end{defn}}
\newenvironment{rem}{\begin{remk}}%
{\end{remk}}
\newenvironment{rems}{\begin{remks}}%
{\hfill$\square$\end{remks}}
\newenvironment{ex}{\begin{exm}}%
{\hfill$\square$\end{exm}}
\newenvironment{exs}{\begin{exms}}%
{\hfill$\square$\end{exms}}
\newenvironment{note}{\begin{notat}}%
{\hfill$\square$\end{notat}}

\newcommand{\thmref}{Theorem~\ref}
\newcommand{\propref}{Proposition~\ref}
\newcommand{\corref}{Corollary~\ref}
\newcommand{\defref}{Definition~\ref}
\newcommand{\lemref}{Lemma~\ref}
\newcommand{\exref}{Example~\ref}
\newcommand{\remref}{Remark~\ref}
\newcommand{\chref}{Chapter~\ref}
\newcommand{\apref}{Appendix~\ref}
\newcommand{\secref}{Section~\ref}
\newcommand{\conjref}{Conjecture~\ref}



\newcommand{\sA}{{\mathcal A}}
\newcommand{\sB}{{\mathcal B}}
\newcommand{\sC}{{\mathcal C}}
\newcommand{\sD}{{\mathcal D}}
\newcommand{\sE}{{\mathcal E}}
\newcommand{\sF}{{\mathcal F}}
\newcommand{\sG}{{\mathcal G}}
\newcommand{\sH}{{\mathcal H}}
\newcommand{\sI}{{\mathcal I}}
\newcommand{\sJ}{{\mathcal J}}
\newcommand{\sK}{{\mathcal K}}
\newcommand{\sL}{{\mathcal L}}
\newcommand{\sM}{{\mathcal M}}
\newcommand{\sN}{{\mathcal N}}
\newcommand{\sO}{{\mathcal O}}
\newcommand{\sP}{{\mathcal P}}
\newcommand{\sQ}{{\mathcal Q}}
\newcommand{\sR}{{\mathcal R}}
\newcommand{\sS}{{\mathcal S}}
\newcommand{\sT}{{\mathcal T}}
\newcommand{\sHH}{{\mathcal H}{\mathcal H}}
\newcommand{\sHC}{{\mathcal H}{\mathcal C}}
\newcommand{\tHH}{\widetilde{{\mathcal H}{\mathcal H}}}
\newcommand{\tHC}{\widetilde{{\mathcal H}{\mathcal C}}}
\newcommand{\tK}{\widetilde{K}}
\newcommand{\sU}{{\mathcal U}}
\newcommand{\sV}{{\mathcal V}}
\newcommand{\sW}{{\mathcal W}}
\newcommand{\sX}{{\mathcal X}}
\newcommand{\sY}{{\mathcal Y}}
\newcommand{\sZ}{{\mathcal Z}}
%\newcommand{\A}{{\mathbb A}}
\newcommand{\B}{{\mathbb B}}
%\newcommand{\C}{{\mathbb C}}
%\newcommand{\D}{{\mathbb D}}
%\newcommand{\E}{{\mathbb E}}
\newcommand{\F}{{\mathbb F}}
%\newcommand{\G}{{\mathbb G}}
\renewcommand{\H}{{\mathbb H}}

\newcommand{\J}{{\mathbb J}}
\newcommand{\M}{{\mathbb M}}
\newcommand{\N}{{\mathbb N}}
\renewcommand{\P}{{\mathbb P}}
\newcommand{\Q}{{\mathbb Q}}
\newcommand{\R}{{\mathbb R}}
\newcommand{\T}{{\mathbb T}}
\newcommand{\U}{{\mathbb U}}
\newcommand{\V}{{\mathbb V}}
\newcommand{\W}{{\mathbb W}}
\newcommand{\X}{{\mathbb X}}
\newcommand{\Y}{{\mathbb Y}}
\newcommand{\Z}{{\mathbb Z}}

\newcommand{\A}{\mathrm{A}}
\newcommand{\C}{\mathrm{C}}
\newcommand{\D}{\mathrm{D}}
\newcommand{\E}{\mathrm{E}}
\newcommand{\G}{\mathrm{G}}

\newcommand{\fa}{{\mathfrak a}}
\newcommand{\fb}{{\mathfrak b}}
\newcommand{\fm}{{\mathfrak m}}
\newcommand{\fp}{{\mathfrak p}}
\newcommand{\fq}{{\mathfrak q}}
\newcommand{\fM}{{\mathfrak M}}
\newcommand{\an}{{\rm an}}
\newcommand{\alg}{{\rm alg}}
\newcommand{\cl}{{\rm cl}}
\newcommand{\Alb}{{\rm Alb}}
\newcommand{\CH}{{\rm CH}}
\newcommand{\mc}{\mathcal}
\newcommand{\mb}{\mathbb}
\newcommand{\surj}{\twoheadrightarrow}
\newcommand{\inj}{\hookrightarrow}
\newcommand{\red}{{\rm red}}
\newcommand{\codim}{{\rm codim}}
\newcommand{\rank}{{\rm rank}}
\newcommand{\Pic}{{\rm Pic}}
\newcommand{\Div}{{\rm Div}}
\newcommand{\Hom}{{\rm Hom}}
%\newcommand{\End}{{\rm End}}
\newcommand{\im}{{\rm im}}
\newcommand{\Spec}{{\rm Spec}}
\newcommand{\sing}{{\rm sing}}
\newcommand{\Char}{{\rm char}}
\newcommand{\Tr}{{\rm Tr}}
\newcommand{\Gal}{{\rm Gal}}
\newcommand{\divf}{{\rm div}}
\newcommand{\Min}{{\rm Min}}
\newcommand{\Max}{{\rm Max \ }}
\newcommand{\supp}{{\rm supp}\,}
\newcommand{\Ker}{{\rm Ker}}
\newcommand{\Img}{{\rm Im}}
\newcommand{\sHom}{{\mathcal{H}{om}}}
\newcommand{\hh}{{\rm ht}}
\newcommand{\sn}{{\rm sn}}
\newcommand{\gp}{{\rm gp}}
\newcommand{\Um}{{\rm Um}}
\newcommand{\Ann}{{\rm Ann}}
\newcommand{\Aut}{{\rm Aut}}
\newcommand{\nil}{{\rm nil}}
\newcommand{\Frac}{{\rm Frac}}
\newcommand{\Jac}{{\rm Jac}}
\newcommand{\sign}{{\rm sign}}

\newcommand{\Nm}{{\operatorname{Nm}}}
\newcommand{\NS}{{\operatorname{NS}}}
\newcommand{\id}{{\operatorname{id}}}
\newcommand{\Zar}{{\text{\rm Zar}}} 
\newcommand{\Ord}{{\mathbf{Ord}}}
\newcommand{\FSimp}{{{\mathcal FS}}}

\newcommand{\eft}{{\operatorname{\mathbf{Eftalg}}}} 
\newcommand{\rft}{{\operatorname{\mathbf{Reftalg}}}} 
\newcommand{\Sch}{{\operatorname{\mathbf{Sch}}}} 
\newcommand{\cosk}{{\operatorname{\rm cosk}}} 
\newcommand{\sk}{{\operatorname{\rm sk}}} 
\newcommand{\subv}{{\operatorname{\rm sub}}}
\newcommand{\bary}{{\operatorname{\rm bary}}}
\newcommand{\Comp}{{\mathbf{SC}}}
\newcommand{\IComp}{{\mathbf{sSC}}}
\newcommand{\Top}{{\mathbf{Top}}}
\newcommand{\holim}{\mathop{{\rm holim}}}
\newcommand{\Holim}{\mathop{{\it holim}}}
\newcommand{\op}{{\text{\rm op}}}
\newcommand{\<}{\langle}
\renewcommand{\>}{\rangle}
\newcommand{\Sets}{{\mathbf{Sets}}}
\newcommand{\del}{\partial}
\newcommand{\fib}{{\operatorname{\rm fib}}}
%\renewcommand{\max}{{\operatorname{\rm max}}}
\newcommand{\bad}{{\operatorname{\rm bad}}}
\newcommand{\Spt}{{\mathbf{Spt}}}
\newcommand{\Spc}{{\mathbf{Spc}}}
\newcommand{\Sm}{{\mathbf{Sm}}}
\newcommand{\SmProj}{{\mathbf{SmProj}}}
\newcommand{\SmAff}{{\mathbf{SmAff}}}
\newcommand{\SmLoc}{{\mathbf{SmLoc}}}
\newcommand{\cofib}{{\operatorname{\rm cofib}}}
\newcommand{\hocolim}{\mathop{{\rm hocolim}}}
\newcommand{\Glu}{{\mathbf{Glu}}}
\newcommand{\can}{{\operatorname{\rm can}}}
\newcommand{\Ho}{{\mathbf{Ho}}}
\newcommand{\GL}{{\operatorname{\rm GL}}}
\newcommand{\SL}{{\operatorname{\rm SL}}}
\newcommand{\El}{{\operatorname{\rm E}}}
\newcommand{\EL}{{\operatorname{\rm EL}}}
\newcommand{\Mat}{{\operatorname{\rm M}}}
\newcommand{\TSp}{{\operatorname{\rm TSp}}}
\newcommand{\ESp}{{\operatorname{\rm ESp}}}
\newcommand{\EGq}{{\operatorname{\rm EGq}}}
\newcommand{\StSp}{{\operatorname{\rm StSp}}}
\newcommand{\Gq}{{\operatorname{\rm Gq}}}
\newcommand{\Sp}{{\operatorname{\rm Sp}}}
\newcommand{\KSp}{{\operatorname{\rm KSp}}}
\newcommand{\K}{{\operatorname{\rm K}}}
\newcommand{\Ss}{{\operatorname{\rm S}}}
\newcommand{\I}{{\operatorname{\rm I}}}
\newcommand{\Id}{{\operatorname{\rm Id}}}
\newcommand{\St}{\mathrm{St}}
\newcommand{\ev}{\mathrm{ev}}

\newcommand{\sq}{\square}
\newcommand{\Ab}{{\mathbf{Ab}}}
\newcommand{\Tot}{{\operatorname{\rm Tot}}}
\newcommand{\loc}{{\operatorname{\rm s.l.}}}
\newcommand{\HZ}{{\operatorname{\sH \Z}}}
\newcommand{\Cyc}{{\operatorname{\rm Cyc}}}
\newcommand{\cyc}{{\operatorname{\rm cyc}}}
\newcommand{\RCyc}{{\operatorname{\widetilde{\rm Cyc}}}}
\newcommand{\Rcyc}{{\operatorname{\widetilde{\rm cyc}}}}
\newcommand{\Sym}{{\operatorname{\rm Sym}}}
\newcommand{\fin}{{\operatorname{\rm fin}}}
\newcommand{\SH}{{\operatorname{\sS\sH}}}
\newcommand{\Wedge}{{\Lambda}}
\newcommand{\eff}{{\operatorname{\rm eff}}}
\newcommand{\rcyc}{{\operatorname{\rm rev}}}
\newcommand{\DM}{{\operatorname{\mathcal{DM}}}}
\newcommand{\GW}{{\operatorname{\rm{GW}}}}
\newcommand{\End}{{\operatorname{\text{End}}}}
\newcommand{\sSets}{{\mathbf{Spc}}}
\newcommand{\Nis}{{\operatorname{Nis}}}
\newcommand{\et}{{\text{\'et}}}
\newcommand{\Cat}{{\mathbf{Cat}}}
\newcommand{\ds}{{/\kern-3pt/}}
\newcommand{\bD}{{\mathbf{D}}}
\newcommand{\bC}{{\mathbf{C}}}
\newcommand{\res}{{\operatorname{res}}}

\renewcommand{\log}{{\operatorname{log}}}
\newcommand{\qf}{{\operatorname{q.fin.}}}
\newcommand{\fil}{\phi}
\newcommand{\barfil}{\sigma}
\newcommand{\Fac}{{\mathop{\rm{Fac}}}}
\newcommand{\Fun}{{\mathbf{Func}}}


\newcommand{\cp}{\coprod}
\newcommand{\ess}{\text{\rm{ess}}}
\newcommand{\dimrel}{\operatorname{\text{dim-rel}}}


\newcommand{\gen}{{\text{gen}}}
\newcommand{\Gr}{{\text{\rm Gr}}}
\newcommand{\Ind}{{\operatorname{Ind}}}
\newcommand{\Supp}{{\operatorname{Supp}}}
\newcommand{\Cone}{{\operatorname{Cone}}}
\newcommand{\Tor}{{\operatorname{Tor}}}
\newcommand{\cok}{{\operatorname{coker}}}
\newcommand{\Br}{{\operatorname{Br}}}
\newcommand{\sm}{{\operatorname{sm}}}
\newcommand{\Proj}{{\operatorname{Proj}}}
\newcommand{\Pos}{{\operatorname{\mathbf{Pos}}}}
\newcommand{\colim}{\mathop{\text{\rm colim}}}
\newcommand{\ps}{{\operatorname{\text{p.s.s.}}}}
\newcommand{\Pro}{{\mathbf{Pro}\hbox{\bf -}}}
\newcommand{\hocofib}{{\operatorname{\rm hocofib}}}
\newcommand{\hofib}{{\operatorname{\rm hofib}}}
\newcommand{\Th}{{\operatorname{\rm Th}}}
\newcommand{\Cube}{{\mathbf{Cube}}}
\newcommand{\Mod}{{\mathbf{Mod}}}
\newcommand{\alt}{{\operatorname{\rm alt}}}
\newcommand{\Alt}{{\operatorname{\rm Alt}}}
\newcommand{\sgn}{{\operatorname{\rm sgn}}}
\newcommand{\dga}{{\rm d.g.a.\!}}
\newcommand{\cdga}{{\rm c.d.g.a.\!}}
\newcommand{\lci}{{\rm l.c.i.\!}}
\newcommand{\sym}{{\operatorname{\rm sym}}}
\newcommand{\DTM}{{\operatorname{DTM}}}
\newcommand{\DCM}{{\operatorname{DCM}}}
\newcommand{\MTM}{{\operatorname{MTM}}}
\newcommand{\Mot}{{\operatorname{Mot}}}
\renewcommand{\sp}{{\widetilde{sp}}}
\newcommand{\NST}{{\mathbf{NST}}}
\newcommand{\PST}{{\mathbf{PST}}}
\newcommand{\Sus}{{\operatorname{Sus}}}
\newcommand{\tr}{{\operatorname{tr}}}
\newcommand{\cone}{{\operatorname{cone}}}
\newcommand{\na}{{\rm{naive}}}
%\newcommand{\ker}{{\rm{Ker}}}
\newcommand{\sem}{{\rm{semi}}}
\newcommand{\Cor}{{\operatorname{Cor}}}
\newcommand{\DMH}{\operatorname{DM}^H}
\newcommand{\gm}{\text{gm}}
\newcommand{\TZ}{{\operatorname{Tz}}}
\renewcommand{\TH}{{\operatorname{TCH}}}
\newcommand{\un}{\underline}
\newcommand{\ov}{\overline}
\renewcommand{\hom}{\text{hom}}
\newcommand{\dgn}{{\operatorname{degn}}}
\renewcommand{\dim}{\text{\rm dim}}
\newcommand{\td}{\mathbf{d}}
\newcommand{\tuborg}{\left\{\begin{array}{ll}}
\newcommand{\sluttuborg}{\end{array}\right.}
\newcommand{\dto}{\dashrightarrow}
\newcommand{\pretr}{{\rm pretr}}
\newcommand{\pr}{{\rm pr}}
\newcommand{\tch}{{\mathcal{TCH}}}
\newcommand{\sfs}{{\rm sfs}}
\newcommand{\fs}{{\rm fs}}
\newcommand{\Sec}{{\rm Sec}}
\newcommand{\zar}{{\rm zar}}
\newcommand{\nis}{{\rm nis}}
\newcommand{\wt}{\widetilde}
\newcommand{\Ass}{{\rm Ass}}
\newcommand{\rU}{{\rm U}}

\newcommand{\bfi}{{\bf{i}}}
\newcommand{\bfj}{{\bf{j}}}
\newcommand{\bfk}{{\bf{k}}}
\newcommand{\bfl}{{\bf{l}}}
\newcommand{\bfm}{{\bf{m}}}
\newcommand{\bfn}{{\bf{n}}}
\newcommand{\bfo}{{\bf{o}}}
\newcommand{\bfp}{{\bf{p}}}
\newcommand{\etl}{{\acute{e}t}}

\def\cO{\mathcal{O}}
\def\cF{\mathcal{F}}
\def\ZDeSXr{\Z(r)^{\mathcal{D}^*}_{S_X}}
\def\ZDeDr{\Z(r)^{\mathcal{D}^*}_{D}}
\def\ZDeXr{\Z(r)^{\mathcal{D}^*}_{X}}
\def\ZDeXrun{\Z(r)^{\mathcal{D}}_{X}}
\def\ZDer#1{\Z(r)^{\mathcal{D}^*}_{#1}}
\def\ZDe#1#2{\Z(#2)^{\mathcal{D}^*}_{#1}}  %usage: $\ZDe{X}{2}$ produces weight 2 modified D-complex on X
\def\ZDeun#1#2{\Z(#2)^{\mathcal{D}}_{#1}}
\def\ol#1{\overline{#1}}
\def\dx{{\rm d}x}



\newcounter{elno}   
\newenvironment{romanlist}{
                         \begin{list}{\roman{elno})
                                     }{\usecounter{elno}}
                      }{
                         \end{list}}
                         
\newcounter{elno-abc}   
\newenvironment{listabc}{
                         \begin{list}{\alph{elno-abc})
                                     }{\usecounter{elno-abc}}
                      }{
                         \end{list}}
\newcounter{elno-abc-prime}   
\newenvironment{listabcprime}{
                         \begin{list}{\alph{elno-abc-prime}')
                   

                  }{\usecounter{elno-abc-prime}}
                      }{
                         \end{list}}






\begin{document}


\title[Graded version of local global principle]{Graded version of local-global principle for Steinberg and Chevalley groups}
\author{Chandan Bhaumik}
%\author{Chandan Bhaumik\textsuperscript{*} and Manoj K. Keshari}


\address{Department of Mathematics, Indian  Institute of Technology Bombay, Powai, Mumbai 400076, India}
 \email{cbhaumik11math@gmail.com, cbhaumik@math.iitb.ac.in}


%\thanks{\textsuperscript{*}Corresponding author. Email: cbhaumik11math@gmail.com (Chandan Bhaumik)}



\dedicatory{}
  

\keywords{Steinberg group, Chevalley group, Local–global principle, Graded ring}        



\subjclass[2020]{20G35, 19C20, 13A02}

\maketitle
\vspace{-.7cm}
% \begin{center}
%    \small Department of Mathematics, IIT Bombay, Powai, Mumbai-400076, India\\ Email: cbhaumik11math@gmail.com
% \end{center}
% \vspace{.5cm}
% \noindent\textit{\textsuperscript{*}Corresponding author. Email: cbhaumik@math.iitb.ac.in}

\begin{quote}\emph{Abstract.} 
In this article, an analog of a graded version of Quillen–Suslin’s local-global principle was established for Steinberg groups $\St(\Phi,-)$ and Chevalley groups $\G(\Phi,-)$ over a reduced irreducible root system $\Phi$.  
\end{quote}
%\end{abstract}
\setcounter{tocdepth}{1} % Ensures subsections appear in the table of contents


%\tableofcontents

\section{Introduction}\label{sec:Intro}
In the context of classical K-theory, the local-global principle was introduced by Quillen \cite{Qu76} and Suslin \cite{Su77} as part of their solution to Serre’s problem on projective modules over polynomial extensions of a field. This principle was initially applied to prove Serre’s conjecture. The solution to Serre’s problem ultimately relied on two key analogs: the $\K_{2}$-analog of Quillen’s local-global principle, which deals with projective modules, and the $\K_{1}$-analog of Suslin's local-global principle, which deals with elementary matrices.\par
Recall that for a commutative ring $R$, denote $\G(\Phi,R)$ as a simply connected Chevalley group (see \cite{St68}) and $\St(\Phi,R)$ as a Steinberg group (see Definition $\ref{def_st_grp}$) with a reduced irreducible root system $\Phi$. The subgroup of $\G(\Phi,R)$ generated by the elementary root unipotents $e_{\alpha}(r)$ for all $\alpha\in\Phi$ and $r\in R$ is called the elementary subgroup of $\G(\Phi,R)$ and is denoted by $\El(\Phi,R)$. There is a well-defined map $\phi:\St(\Phi,R)\to\G(\Phi,R)$ sending $x_{\alpha}(r)\mapsto e_{\alpha}(r)$, and the image of $\phi$ is $\El(\Phi,R)$, which is a normal subgroup of $\G(\Phi,R)$ when $\rank(\Phi)\geq 2$. The groups $\K_{1}(\Phi,R)$ and $\K_{2}(\Phi,R)$ are defined as the cokernel and kernel of $\phi$, respectively
\begin{center}
        \begin{tikzcd}
1 \arrow[r] & \K_{2}(\Phi,R) \arrow[r]           & \St(\Phi, R) \arrow[r,"\phi"]                 & \G(\Phi, R) \arrow[r]  & \K_{1}(\Phi,R) \arrow[r]               & 1. 
\end{tikzcd}
\end{center}
The $\K_{1}$-analog of the local-global principle was established for Chevalley groups of type $\A_{n}$ by Suslin \cite{Su77}. Subsequent work extended it to type $\D_{n}$ by Suslin–Kopeiko \cite{SuKo82}, to type $\C_{n}$ by Grunewald–Mennicke–Vaserstein \cite{GMV91}, and to all  Chevalley groups by Abe \cite{Ab83} and Wendt \cite{We10}. Turning to the $\K_{2}$-analog of the local-global principle, it was proven for Chevalley groups of type $\A_{n}$ by Tulenbaev \cite{Tu83}. Recently, it was extended to type $\E_{n}$ by Sinchuk \cite{Si16}, to type $\D_{n}$ by Lavrenov-Sinchuk \cite{LaSi17}, and to type $\C_{n}$ by Lavrenov \cite{La18}. The present article aims to revisit these results in the setting of commutative $\N$-graded rings.\par
The graded version of the $\K_{1}$-analog of the local-global principle has been studied for classical groups in \cite{BaSi20,Gu92}. Meanwhile, for Chevalley groups and Steinberg groups, the graded versions of both the $\K_{1}$-analog and $\K_{2}$-analog are expected to hold. We therefore establish the graded versions of the local-global principle for Steinberg groups and Chevalley groups. These results constitute the main theorems of the present article and are stated below.
\begin{thm}[Theorem \ref{thm_lgp_stgp}]\label{intro_thm_lgp_stgp}
    Let $R=\bigoplus_{i\geq 0}R_{i}$ be a graded ring and $\Phi_{n}$ be a either root system of type simply-laced or $\C_{n}$ of rank $\geq 3$. Suppose $\theta\in \St(\Phi_{n},R, R^{+})$. Then \begin{enumerate}
        \item $\theta=1$ if and only if $\lambda^{\ast}_{\mathfrak{m}}(\theta)=1\in \St(\Phi_{n},R_{\mathfrak{m}})$ for all maximal ideals $\mathfrak{m}$ of $R_{0}$.
        \item $\theta\in\Img(l_{R})$ if and only if $\lambda^{\ast}_{\mathfrak{m}}(\theta)\in \Img(l_{R_{\mathfrak{m}}})$ for all maximal ideals $\mathfrak{m}$ of $R_{0}$.
    \end{enumerate}
\end{thm}
\begin{thm}[Theorem \ref{thm_lgp_chgp}]\label{intro_thm_lgp_chgp}
     Let $R=\bigoplus_{i\geq 0}R_{i}$ be a graded ring and $\Phi$ be a root system of rank $\geq 2$. Assume that $R$ has no residue fields of $2$ elements if $\Phi=\C_{2}~\text{or}~\G_{2}$. Suppose $\theta\in \G(\Phi,R,R^{+})$. Then $\theta\in \El(\Phi,R)$ if and only if $\lambda^{\ast}_{\mathfrak{m}}(\theta)\in \El(\Phi,R_{\mathfrak{m}})$ for all maximal ideals $\mathfrak{m}$ of $R_{0}$.
\end{thm}
The centrality of $\K_{2}$ (see \cite{LaSi17}) follows from the $\K_{2}$-analog of Quillen’s local-global principle, just as the normality theorem for the elementary subgroup is derived from the $\K_{1}$-analog of Suslin's local-global principle. In other words, the local-global principle plays a foundational role in establishing both the normality of $\El(\Phi,R)$ and the centrality of $\St(\Phi,R)$. To extend these results from the polynomial ring to the graded ring,  one must employ a combination of three key techniques: the methodology of Quillen-Suslin’s argument, the dilation principle, and Swan--Weibel’s homotopy technique, which allows one to connect polynomial rings to graded rings via a homotopy map.
\section{Preliminaries and Notations}\label{sec_preli}
{\bf Throughout the article, all rings are assumed to be commutative with unity, and all root systems are reduced and irreducible of rank $\geq 2$ unless otherwise stated.} 
\begin{defn}\label{def_st_grp}
Let $A$ be a ring and $\Phi$ be a root system. The {\it Steinberg group} $\St(\Phi, A)$ is the group generated by symbols $x_\alpha(u)$, where $\alpha \in \Phi$ and $u \in A$, subject to the following relations:
\begin{align}
    &    x_\alpha(u) \, x_\alpha(v) = x_\alpha(u+v), \qquad u, v \in R,~ \alpha\in\Phi;\tag{St1} \label{St1}\\[4pt]
    &    [x_\alpha(u), x_\beta(v)] = \prod_{\substack{i\alpha + j\beta \in \Phi}} x_{i\alpha + j\beta}\bigl(C_{i,j,\alpha,\beta}\, u^i v^j\bigr), \qquad u, v \in R,~\alpha, \beta \in \Phi, ~\alpha\neq -\beta,\tag{St2} \label{St2}
\end{align}
where $[p,q] = pqp^{-1}q^{-1}$, the product is taken over all positive integers $i,j$ such that $i\alpha + j\beta \in \Phi$, and $C_{i,j,\alpha,\beta} \in \mathbb{Z}$ are the Chevalley structure constants (depending only on $\Phi$). More information on Chevalley groups, the structure constants, and various types of root systems $\Phi$ can be found in \cite{St68,VaPl96}. 
\end{defn}
\begin{defn}
    A root system $\Phi$ is called {\it simply-laced} if the Dynkin diagram of $\Phi$ does not contain double bonds (i.e., $\Phi$ has type $\A_n,\D_n$ and $\E_{6,7,8}$). 
\end{defn}
For an ideal $I$ of $A$, the subgroup $\St(\Phi,I)$ is spanned by symbols $x_\alpha(s)$, where $\alpha \in \Phi$ and $s \in I$. An ideal $I$ is called a {\it splitting ideal} if the canonical projection ring homomorphism $\pi:A \twoheadrightarrow A/I$ splits, i.e., there exists $\sigma:A/I\to A$ such that $\pi\sigma=1$. It is noted that $TA[T]$ is a splitting ideal of $A[T]$.
\begin{defn}
  For a splitting ideal $I$, the {\it relative Steinberg group} $\St(\Phi,A,I)$ is defined as the kernel  $\Ker(\pi^{\ast}:\St(\Phi,A)\twoheadrightarrow\St(\Phi,A/I))$. Obviously, $\St(\Phi,A,I)$ is coincides with the normal closure of $\St(\Phi,I)$ in $\St(\Phi,A)$.  
\end{defn}
In the case of an arbitrary ideal $I$, the relative Steinberg group $\St(\Phi,A,I)$  can be defined as a central extension of $\Ker(\pi^{\ast})$. For the purposes of the present text, we concentrate on the relative Steinberg groups corresponding to splitting ideals. For more information regarding the general case, we refer the reader to \cite{Si16}.
\begin{defn}(Universal Property)
 Let $\mathfrak{F}(X)$ denote the free group on the generating set $X$. If $G$ is any group with a generating set $X$ and certain defining relations $u_i=v_i$, there is an isomorphism $\mathfrak{F}(X)/\mathfrak{D}\cong G$ taking $\bar{x}\mapsto x$ for each $x\in X$, where $\mathfrak{D}$ is a normal subgroup of $\mathfrak{F}(X)$ generated by the
conjugates of the relators $u_iv_i^{-1}$. Then $\mathfrak{F}(X)$ has the following universal property: If $H$ is any group and each function $f: X\to H$ extends uniquely to a group homomorphism $\varphi:\mathfrak{F}(X)\to H$ that makes the diagram commute.
\begin{center}
    \begin{tikzcd}
X \arrow[rd, "f"] \arrow[d]      &   \\
\mathfrak{F}(X) \arrow[d] \arrow[r, "\varphi"] & H \\
\mathfrak{F}(X)/\mathfrak{D} \arrow[ru, "f^{*}" ']            &  
\end{tikzcd}
\end{center}\label{defn_uniprop}
\end{defn}
Now, to define a group homomorphism $f^{*}:G\to H$, it suffices to choose a function $f: X\to H$ for
which, when $x$ is replaced by $f(x)$ and $x^{-1}$ by $f(x)^{-1}$, the defining relations $u_i=v_i$ for $G$ become true equations in $H$. Indeed, the homomorphism $\varphi$ extending $f$ takes $\mathfrak{D}$ to $1_{H}$; thus, it induces a homomorphism $f^{*}$ taking $\bar{x}$ to $f(x)$ for $x\in X$. Thus, the composite map $f^{*}:\mathfrak{F}(X)/\mathfrak{D}\to H$ is a homomorphism taking $x\in X$ to $f(x)\in H$. For more details on universal property and classical $\K$-theory, we refer the reader to \cite{HaMe89}.

\begin{lem}
    If $f: A\to B$  is a ring homomorphism, then $f$ induces a group homomorphism $f^{\ast}:\St(\Phi,A)\to \St(\Phi, B)$.\label{lem_st_hom}
\end{lem}
\begin{proof}
    Since any element of $\St(\Phi,A)$ can be written as a finite product of elements of the form $x_{\alpha}(u)$, $\alpha\in \Phi,~u\in A$, it is enough to define a map on generators. We define a map $f^{\ast}(x_{\alpha}(u))=x_{\alpha}(f(u))$ where $\alpha\in \Phi$. Due to the universal property $\ref{defn_uniprop}$, we need to verify that this map preserves the Steinberg relations, i.e., replacing $u$ by $f(u)$ and $v$ by $f(v)$ in the defining relations {\rm (St1)} and {\rm (St2)} for $\St(\Phi,A)$ yields equations in $\St(\Phi,B)$. However, this is indeed true because $f$ is a ring homomorphism.
\end{proof}
One of the key components for handling the graded situation is the well-known ``Swan--Weibel's homotopy trick".
Let $R=R_{0}\oplus R_{1}\oplus R_{2}\oplus\cdots$ be a graded ring and $R^{+}=R_{1}\oplus R_{2}\oplus\cdots$ denote the irrelevant ideal of $R$. An element $r\in R$ will be denoted by $r=r_0+r_1+r_2+\cdots$ where $r_i\in R_i$ for each $i$, and all but finitely many $r_{i}$ are zero.
\begin{defn}\label{def_sw_map}
    For a graded ring $R=\bigoplus_{i\geq 0}R_{i}$, the {\it Swan--Weibel homotopy map} $h:R\to R[T]$ is defined as \[h(r_0+r_1+\cdots+r_m)=r_0+r_1 T+\cdots r_m T^{m}, ~(r_{i}\in R_{i}).\]
For any $R_0$-algebra $A$ and any $a\in A$, consider an evaluation map $\ev_{a}:R[T]\to A$, which is the $R_0$-algebra homomorphism sending $T\mapsto a$.
\end{defn}
We denote the image $r$ in $A$ by $r^{h}(a):=(\ev_{a}\circ h)(r)$. If $A=R$, then it is easy to observe that $\ev_{0}\circ h: R\to R$ is precisely the projection onto the degree-zero part and $\ev_{1}\circ h: R\to R$ is precisely the identity map. For a fixed element $c\in R_0$, similarly denote $r^{h}(c):=h(r)(c)= (\ev_{c}\circ h)(r)$, and note that $r^{h}(0)=r_0$ and $r^{h}(1)=r$. In order to emphasize the choice of $A$, we denote by $\ev\left [\frac{R[T]\to A}{T\mapsto a}\right ]$ for $R_0$-algebra homomorphism $\ev_{a}$. 
\begin{rem}\label{rem_st_ev}
In view of Lemma $\ref{lem_st_hom}$, the map $h:R\to R[T]$ (resp. $\ev_{a}:R[T]\to A$) induces a group homomorphism $h^{\ast}:\St(\Phi,R)\to \St(\Phi,R[T])$ (resp. $\ev_{a}^{\ast}:\St(\Phi,R[T])\to \St(\Phi,A)$) that sends the generators $x_{\alpha}(r)$ to $x_{\alpha}(h(r))$ (resp. $x_{\alpha}(p(T))$ to $x_{\alpha}(p(a))$) where $\alpha\in \Phi$. For $\theta\in \St(\Phi,R)$ and $c\in R_{0}$, analogously, $\theta^{h^{\ast}}(c):=(\ev_{c}^{\ast}\circ h^{\ast})(\theta)$ is defined. For more details on Swan--Weibel's homotopy trick, we refer the reader to \cite[Chapter V]{Lambook}.
\end{rem}
\textbf{Notations:} For a ring $A$, we denote by $\lambda_{a}:A\to A_{a}$ and $\lambda_{\mathfrak{m}}:A\to A_{\mathfrak{m}}$ the homomorphisms of localization at an element $a\in A$ and at a maximal ideal $\mathfrak{m}$ of $A$ respectively. If $f:A\to B$ is a ring homomorphism, its corresponding Steinberg group homomorphism is denoted by $f^{\ast}:\St(\Phi,A)\to \St(\Phi,B)$. If $\Phi_{n}$ is an irreducible root system of rank $n$, we denote by $l_{A}:\St(\Phi_{n-1},A)\to \St(\Phi_{n},A)$ the homomorphism of Steinberg groups. For visual understanding, it is noted that we have the following commutative diagram
\begin{center}
    \begin{tikzcd}
\St(\Phi_{n-1},A) \arrow[r, "l_{A}"] \arrow[d,"f^{\ast}"] & \St(\Phi_{n},A) \arrow[d, "f^{\ast}"] \\
\St(\Phi_{n-1},B) \arrow[r, "l_{B}"]                & \St(\Phi_{n},B).               
\end{tikzcd}
\end{center}
\section{Local-Global Principle for Stienberg Groups}\label{sec:lgp}
In this section, we assume that $\Phi_{n}$ is a root system either of type simply-laced (i.e., $\A_{n\geq 3},\D_{n\geq 3}~\text{or}~\E_{n=6,7,8}$) or $\C_{n\geq 3}$. We begin with the following observation.
\begin{lem}
     Let $R=\bigoplus_{i\geq 0}R_{i}$ be a graded ring and `overbar' denote the reduction modulo the ideal $R^{+}$. For $\theta\in \St(\Phi,R)$ one has $\theta^{h^{\ast}}(0)=\overline{\theta}$ and $\theta^{h^{\ast}}(1)=\theta$. \label{lem_st_hom_sw}
\end{lem}
\begin{proof}
    An element $\theta\in \St(\Phi,R)$ has the form $\prod_{i=1}^{k}x_{\alpha_{i}}(r_{i})$, where $\alpha_{i}\in\Phi$ and $r_{i}=\oplus_{j=0}^{m_{i}} r_{ij}$ with $r_{ij}\in R_{j}$. Then 
    \[\begin{split}
        h^{\ast}(\theta) &=\prod_{i=1}^{k}h^{\ast}(x_{\alpha_{i}}(r_{i}))\\&=\prod_{i=1}^{k}x_{\alpha_{i}}(h(\oplus_{j=0}^{m_{i}} r_{ij}))~(\text{by Remark}~ \ref{rem_st_ev})\\
        &=\prod_{i=1}^{k}x_{\alpha_{i}}(\oplus_{j=0}^{m_{i}} r_{ij}T^{j}).
    \end{split}\]
    Now, it is easy to observe that $\theta^{h^{\ast}}(0)=(\ev_{0}^{\ast}\circ h^{\ast})(\theta)=\prod_{i=1}^{k}x_{\alpha_{i}}(r_{i0})=\overline{\theta}$ and $\theta^{h^{\ast}}(1)=(\ev_{1}^{\ast}\circ h^{\ast})(\theta)=\theta$.
\end{proof}
The following lemma is known as the dilation principle for Steinberg groups.
\begin{lem}
    Let $A$ be an arbitrary ring and $a\in A$. Suppose $\theta(Z)\in \St(\Phi_{n},A[Z], ZA[Z])$.\begin{enumerate}
        \item If $\lambda_{a}^{\ast}(\theta(Z))=1$, then for sufficiently large $k$ one has $\theta(a^kZ)=1$.
        \item If $\lambda_{a}^{\ast}(\theta(Z))\in\Img(l_{A_{a}[Z]})$, then for sufficiently large $k$ one has $\theta(a^kZ)\in\Img(l_{A[Z]})$.
    \end{enumerate} \label{lem_dilation}
\end{lem}
\begin{proof}
    $(1)$ \cite[Proposition 2.2]{LaSi17} for $\Phi_n=\A_n,\D_n,\E_{6,7,8}$ and \cite[Lemma 3.1]{La18} for $\Phi_n=\C_n$. $(2)$ \cite[Lemma 5.1]{Si25} for $\Phi_n=\A_n,\D_n,\E_{6,7,8}$ and \cite[Lemma 3.1]{La18} for $\Phi_n=\C_n$.
\end{proof}
The following lemma
is a crucial step towards the proof of Theorem $\ref{intro_thm_lgp_stgp}$.
\begin{lem}
    Let $R=\bigoplus_{i\geq 0}R_{i}$ be a graded ring and $a,b\in R_{0}$ be two co-maximal elements. Suppose $\theta\in \St(\Phi_{n},R, R^{+})$. 
     \begin{enumerate}
        \item If $\lambda^{\ast}_{a}(\theta)$ and $\lambda^{\ast}_{b}(\theta)$ are equal to $1$, then $\theta=1$.
        \item If $\lambda^{\ast}_{a}(\theta)\in\Img(l_{R_{a}})$ and $\lambda^{\ast}_{b}(\theta)\in\Img(l_{R_{b}})$, then $\theta\in\Img(l_{R})$. 
    \end{enumerate}\label{lem_comax}
\end{lem}
\begin{proof}
     Consider an element $\xi:=\xi(X,Y)=\theta^{h^{\ast}}(X)\cdot(\theta^{h^{\ast}}(X+Y))^{-1}\in\St(\Phi_n,A[Y])$, where $A=R[X]$. This means that \[\xi(X,Y)=\left(\ev\left [\frac{R[T]\to A[Y]}{T\mapsto X}\right ]^{\ast}\circ h^{\ast}\right)(\theta)\cdot\left(\ev\left [\frac{R[T]\to A[Y]}{T\mapsto X+Y}\right ]^{\ast}\circ h^{\ast}\right)(\theta^{-1}).\] Now, it is easy to observe that $\xi(X,0)=1$. Hence, $\xi(X,Y)$ belongs to the kernel of $\St(\Phi_{n},A[Y])\to\St(\Phi_{n},A)$. Thus, we can assume that $\xi(X,Y)\in\St(\Phi_{n},A[Y], YA[Y])$. The image of $\xi(X,Y)$ in the ring $\St(\Phi_n,A_{a}[Y])$ is \[\begin{split}
         \lambda^{\ast}_{a}(\xi)&=\left(\lambda^{\ast}_{a}\left(\ev\left [\frac{R[T]\to A[Y]}{T\mapsto X}\right ]^{\ast}\circ h^{\ast}\right)\right)(\theta)\cdot\left(\lambda^{\ast}_{a}\left(\ev\left [\frac{R[T]\to A[Y]}{T\mapsto X+Y}\right ]^{\ast}\circ h^{\ast}\right)\right)(\theta^{-1})\\
         &=\left(\ev\left [\frac{R_{a}[T]\to A_{a}[Y]}{T\mapsto X}\right ]^{\ast}\circ h_{a}^{\ast}\right)( \lambda^{\ast}_{a}(\theta))\cdot\left(\ev\left [\frac{R_{a}[T]\to A_{a}[Y]}{T\mapsto X+Y}\right ]^{\ast}\circ h_{a}^{\ast}\right)( \lambda^{\ast}_{a}(\theta^{-1}))\\
         &=(\lambda^{\ast}_{a}(\theta))^{h_{a}^{\ast}}(X)\cdot((\lambda^{\ast}_{a}(\theta))^{h_{a}^{\ast}}(X+Y))^{-1}.
     \end{split}\] \\
     $(1)$ Since $\lambda^{\ast}_{a}(\xi)=1\in\St(\Phi_n,A_{a}[Y])$ and analogously $\lambda^{\ast}_{b}(\xi)=1\in\St(\Phi_{n},A_{b}[Y])$, by Lemma $\ref{lem_dilation}(1)$, for some natural number $k$, we have $\xi(X,a^{k}Y)=1$ and $\xi(X,b^{k}Y)=1$. We choose $u,v\in R_{0}$ such that $ua^{k}+vb^{k}=1$. This implies $\xi(X,-ua^{k}Y)=1$ and $\xi(X,-vb^{k}Y)=1$. From Lemma $\ref{lem_st_hom_sw}$, we have $\theta^{h^{\ast}}(0)=1$. Then
     \[\begin{split}
        \theta &=\theta^{h^{\ast}}(1)(\theta^{h^{\ast}}(0))^{-1}\\&=\theta^{h^{\ast}}(1)(\theta^{h^{\ast}}(ua^{k}))^{-1}\theta^{h^{\ast}}(ua^{k})(\theta^{h^{\ast}}(0))^{-1}\\
        &=\xi(1,-vb^{k})\xi(ua^{k},-ua^{k})\\
        &=1.
    \end{split}\]
    $(2)$ The hypothesis implies that $\lambda^{\ast}_{a}(\xi)\in\Img(l_{A_{a}[Y]})$ and $\lambda^{\ast}_{b}(\xi)\in\Img(l_{A_{b}[Y]})$. By Lemma $\ref{lem_dilation}(2)$, both $\xi(X,-ua^{k}Y)$ and  $\xi(X,-vb^{k}Y)$ lie in the image of $l_{A[Y]}$ as in the case $(1)$. Then the identities $\theta=\xi(1,-vb^{k})\xi(ua^{k},-ua^{k})$ implies that $\theta\in\Img(l_{R})$.
\end{proof}
We now prove Theorem $\ref{intro_thm_lgp_stgp}$, which we restate here for the reader’s convenience.
\begin{thm}\label{thm_lgp_stgp}
    Let $R=\bigoplus_{i\geq 0}R_{i}$ be a graded ring. Suppose $\theta\in \St(\Phi_{n},R, R^{+})$. Then \begin{enumerate}
        \item $\theta=1$ if and only if $\lambda^{\ast}_{\mathfrak{m}}(\theta)=1\in \St(\Phi_{n},R_{\mathfrak{m}})$ for all maximal ideals $\mathfrak{m}$ of $R_{0}$.
        \item $\theta\in\Img(l_{R})$ if and only if $\lambda^{\ast}_{\mathfrak{m}}(\theta)\in \Img(l_{R_{\mathfrak{m}}})$ for all maximal ideals $\mathfrak{m}$ of $R_{0}$.
    \end{enumerate}
\end{thm}
\begin{proof}
    We first prove that the following set is an ideal of $R_{0}$ (called the Quillen ideal). \[\mathfrak{Q}(R_{0})=\begin{cases}
        \{a\in R_{0}~\mid~\lambda^{\ast}_{a}(\theta)=1\}  & \text{for (1),} \\
    \{a\in R_{0}~\mid~\lambda^{\ast}_{a}(\theta)\in \Img(l_{R_{a}})\} & \text{for (2)}.
    \end{cases}\]
    Clearly $0\in\mathfrak{Q}(R_{0})$. If $a,b\in \mathfrak{Q}(R_{0})$ and $c\in aR_{0}+bR_{0}$, then clearly $ac,bc\in \mathfrak{Q}(R_{0})$. Note that $a$ and $b$ are comaximal elements of $(R_{0})_{c}$. Since $\lambda^{\ast}_{c}(\theta)\in\St(\Phi_{n},R_{c},R^{+}_{c})$, apply Lemma $\ref{lem_comax}$ to the ring $R_{c}$, we get $\lambda^{\ast}_{c}(\theta)=1$ (resp. $\lambda^{\ast}_{c}(\theta)\in\Img(l_{R_{c}})$), and hence $c\in\mathfrak{Q}(R_{0})$. Therefore, $\mathfrak{Q}(R_{0})$ is an ideal of $R_{0}$. \par
    It suffices to show if part of the statement. To get this we need to prove that $\mathfrak{Q}(R_{0})=R_{0}$, i.e., $1\in \mathfrak{Q}(R_{0})$. If otherwise, there exists a maximal ideal $\mathfrak{m}$ of $R_{0}$ such that $\mathfrak{Q}(R_{0})\subseteq \mathfrak{m}$. Since $\lambda^{\ast}_{\mathfrak{m}}(\theta)=1$ (resp. $\lambda^{\ast}_{\mathfrak{m}}(\theta)\in \Img(l_{R_{\mathfrak{m}}})$), there exists $a\in R_{0}\setminus\mathfrak{m}$ such that $\lambda^{\ast}_{a}(\theta)=1$ (resp. $\lambda^{\ast}_{\mathfrak{m}}(\theta)\in \Img(l_{R_{a}})$). Then $a\in\mathfrak{Q}(R_{0})\subseteq \mathfrak{m}$ which is a contradiction. 
    %But this is indeed true according to the hypothesis.
\end{proof}
%\section{Auxiliary Result for Chevalley Group}\label{sec_split}
\section{Local-Global Principle for Chevalley Groups}
In this section, we deduce an analog of the local-global principle for the Chevalley group over graded rings, and we assume that $R$ has no residue fields of $2$ elements when $\Phi=\C_{2}~\text{or}~\G_{2}$. We first recall the definition of a relative elementary subgroup.
\begin{defn}
     For an ideal $I$ of a ring $A$, the congruence subgroup $\G(\Phi,A,I)$ is the kernel of the canonical map $\G(\Phi,A)\to \G(\Phi,A/I)$. The subgroup $\El(\Phi,I)$ of $\G(\Phi,A,I)$ is spanned by elementary root unipotents $e_{\alpha}(s)$ where $s\in I$ and the relative elementary subgroup $\El(\Phi,A,I)$ is a normal closure of $\El(\Phi,I)$ in $\El(\Phi,A)$.
\end{defn}    
Let $R=\bigoplus_{i\geq 0}R_{i}$ be a graded ring. The Swan-Weibel’s homotopy map $h:R\to R[T]$ induced a group homomorphism  $h^{\ast}:\G(\Phi,R)\to \G(\Phi,R[T])$. As in Remark $\ref{rem_st_ev}$, if $\theta\in \G(\Phi,R)$, for $c\in R_{0}$, we denote by $\theta^{h^{\ast}}(c):=(\ev_{c}^{\ast}\circ h^{\ast})(\theta)$.\par
The following lemma is similar to Lemma $\ref{lem_st_hom}$. Hence, we skip the proof.
\begin{lem}
     If $f: A\to B$  is a ring homomorphism, then $f$ induces a group homomorphism $f^{\ast}:\El(\Phi,A)\to \El(\Phi, B)$ sending the generators $e_{\alpha}(a)$ to $e_{\alpha}(f(a))$.\label{lem_el_hom}
\end{lem}
The following lemma generalizes the result of Suslin's \cite[Lemma 3.4]{Su77}, where he proved it over a polynomial ring for an elementary linear group.
\begin{lem}\label{lem_cong_chgp}
    Let $R=\bigoplus_{i\geq 0}R_{i}$ be a graded ring. Suppose $\theta\in \G(\Phi,R)$, and $s\in R_{0}$. If $\lambda^{\ast}_{s}(\theta)\in\El(\Phi,R_{s})$, then there exists an integer $k$ such that, for two indeterminates $X,Y$, we have \[\theta^{h^{\ast}}(X+s^{k}Y)\cdot\theta^{h^{\ast}}(X)^{-1}\in\El(\Phi,R[X,Y],(Y)).\] In particular, for any $c,d\in R_{0}$ with $c\equiv d ~(\text{mod} ~s^{k}R_{0})$, we have $\theta^{h^{\ast}}(c)\cdot\theta^{h^{\ast}}(d)^{-1}\in\El(\Phi,R)$.
\end{lem}
\begin{proof}
   Denote $\Gamma:=\Gamma(Z)=\theta^{h^{\ast}}(Z)=\left(\ev\left [\frac{R[T]\to R[Z]}{T\mapsto Z}\right ]^{\ast}\circ h^{\ast}\right)(\theta)\in \G(\Phi,R[Z])$. Since $\lambda^{\ast}_{s}(\theta)\in\El(\Phi,R_{s})$, by Lemma $\ref{lem_el_hom}$, we have $\lambda^{\ast}_{s}(\Gamma)=\left(\ev\left [\frac{R_s[T]\to R_s[Z]}{T\mapsto Z}\right ]^{\ast}\circ h_s^{\ast}\right)(\lambda^{\ast}_{s}(\theta))\in\El(\Phi,R_{s}[T])$. Consider the following element in $\G(\Phi,R[X,Y,Z],(Y))$, 
    \[\begin{split}
    \Gamma'(X,Y,Z)&=\Gamma((X+Y)Z)\cdot\Gamma(XZ)^{-1}\\&=\theta^{h^{\ast}}((X+Y)Z)\cdot\theta^{h^{\ast}}(XZ)^{-1}\\&=\ev\left [\frac{R[T]\to R[X,Y,Z]}{T\mapsto (X+Y)Z}\right ]^{\ast}(h^{\ast}(\theta))\cdot\ev\left [\frac{R[T]\to R[X,Y,Z]}{T\mapsto XZ}\right ]^{\ast}(h^{\ast}(\theta^{-1})).
     \end{split}\] 
     Since $\lambda^{\ast}_{s}(\Gamma)\in\El(\Phi,R_{s}[Z])$, we have $\lambda^{\ast}_{s}(\Gamma')\in\El(\Phi,R_{s}[X,Y,Z])$. Applying \cite[Lemma 1.4]{ApSt13} to the ring $R_{s}[X,Y,Z]$ and the splitting ideal $YR_{s}[X,Y,Z]$, we have \[\G(\Phi,R_{s}[X,Y,Z],(Y))\cap\El(\Phi,R_{s}[X,Y,Z])=\El(\Phi,R_{s}[X,Y,Z],(Y)).\] Therefore, $\lambda^{\ast}_{s}(\Gamma')\in\El(\Phi,R_{s}[X,Y,Z],(Y))$. Then, by \cite[Theorem 6.3]{ApSt13}, there exists an integer $k$ such that $\Gamma'(X,s^{k}Y,Z)\in\El(\Phi,R[X,Y,Z],(Y))$. Hence, substituting $Z=1$, we have that \[\Gamma'(X,s^{k}Y):=\theta^{h^{\ast}}(X+s^{k}Y)\cdot\theta^{h^{\ast}}(X)^{-1}\in\El(\Phi,R[X,Y],(Y)).\] 
   For any $c,d\in R_{0}$ with $c=d+s^{k}b$ for some $b\in R_{0}$ and specializing $X=d$ and $Y=b$, we get \[\Gamma'(d,s^{k}b):=\theta^{h^{\ast}}(c)\cdot\theta^{h^{\ast}}(d)^{-1}\in\El(\Phi,R).\] 
   This gives the desired conclusion.
\end{proof}
\begin{thm}\label{thm_lgp_chgp}
     Let $R=\bigoplus_{i\geq 0}R_{i}$ be a graded ring. Suppose $\theta\in \G(\Phi,R)$ with $\theta^{h^{\ast}}(0)=1$. Then $\theta\in \El(\Phi,R)$ if and only if $\lambda^{\ast}_{\mathfrak{m}}(\theta)\in \El(\Phi,R_{\mathfrak{m}})$ for all maximal ideals $\mathfrak{m}$ of $R_{0}$.
\end{thm}
\begin{proof}
It suffices to show if part of the statement. We shall first prove that if $\lambda^{\ast}_{a}(\theta)\in\El(\Phi,R_{a})$ and $\lambda^{\ast}_{b}(\theta)\in\El(\Phi,R_{b})$ with $a,b$ are two co-maximal elements of $R_{0}$, then $\theta\in\El(\Phi,R)$.\par
By Lemma $\ref{lem_cong_chgp}$, we can choose $k$ such that $\theta^{h^{\ast}}(c)\theta^{h^{\ast}}(d)^{-1}\in\El(\Phi,R)$ whenever $c\equiv d~(\text{mod} ~a^{k}R_{0})$ or $c\equiv d~(\text{mod} ~b^{k}R_{0})$. Note that $a^{k}u+b^{k}v=1$ for some $u,v\in R_{0}$. Since $\theta^{h^{\ast}}(0)=1$, we write 
\[\begin{split}
        \theta &=\theta^{h^{\ast}}(1)(\theta^{h^{\ast}}(ua^{k}))^{-1}\theta^{h^{\ast}}(ua^{k})(\theta^{h^{\ast}}(0))^{-1}.
    \end{split}\]
 Then $\theta\in\El(\Phi,R)$ since $a^{k}u\equiv 1 ~(\text{mod} ~b^{k}R_{0})$ and $a^{k}u\equiv 0~(\text{mod} ~a^{k}R_{0})$.\par
As in Theorem $\ref{thm_lgp_stgp}$, it is easy to verify that the set $\mathfrak{Q}(\theta)=\{a\in R_{0}~\mid~\lambda^{\ast}_{a}(\theta)\in\El(\Phi,R_{a})\}$ is an ideal of $R_{0}$. To show $\theta\in \El(\Phi,R)$ we need to prove that $1\in \mathfrak{Q}(R_{0})$. If otherwise, there exists a maximal ideal $\mathfrak{m}$ such that $\mathfrak{Q}(R_{0})\subseteq \mathfrak{m}$. Since $\lambda^{\ast}_{\mathfrak{m}}(\theta)\in \El(\Phi,R_{\mathfrak{m}})$, there exists $a\in R_{0}\setminus\mathfrak{m}$ such that $\lambda^{\ast}_{a}(\theta)\in\El(\Phi,R_{a})$. Then $a\in\mathfrak{Q}(R_{0})\subseteq \mathfrak{m}$ which is a contradiction.
\end{proof}
\begin{cor}
  Let $R=\bigoplus_{i\geq 0}R_{i}$ be a graded ring and $\theta\in \G(\Phi,R)$. If $\lambda^{\ast}_{\mathfrak{m}}(\theta)\in \G(\Phi,(R_{0})_{\mathfrak{m}})\cdot\El(\Phi,R_{\mathfrak{m}})$ for all maximal ideals $\mathfrak{m}$ of $R_{0}$, then $\theta\in \G(\Phi,R_{0})\cdot\El(\Phi,R)$.
\end{cor}
\begin{proof}
   Let `overbar' denote the reduction modulo the ideal $R^{+}$. After left-multiplying $\theta$ by $\overline{\theta}^{-1}$, we can assume that $\theta^{h^{\ast}}(0)=1$. For all maximal ideals $\mathfrak{m}$ of $R_{0}$, we have the decomposition $\lambda^{\ast}_{\mathfrak{m}}(\theta)=\eta\cdot \gamma$ where $\eta\in \G(\Phi,(R_{0})_{\mathfrak{m}})$ and $\gamma\in \El(\Phi,R_{\mathfrak{m}})$. Now, reduction modulo $R^{+}$, we get $I=\eta\cdot\overline{\gamma}$, and hence $\eta=\overline{\gamma}^{-1}$. Thus $$\lambda^{\ast}_{\mathfrak{m}}(\theta)=\overline{\gamma}^{-1}\cdot\gamma\in\El(\Phi,R_{\mathfrak{m}}).$$ Then, by Theorem $\ref{thm_lgp_chgp}$, we get $\theta\in\El(\Phi,R)$. Therefore $\theta\in \G(\Phi,R_{0})\cdot\El(\Phi,R)$ since the original $\theta$ has already been modified by a left factor from $\G(\Phi,R_{0})$. %see lam book page 207
\end{proof}
 
\textbf{Acknowledgments:} The author thanks Prof. Manoj Kumar Keshari for suggesting the problem. This material is based on work supported by IIT Bombay Postdoctoral Research Fellowship.
\iffalse 
\\ \\
\textbf{\Large Declarations} \\\\
\textbf{\large Funding} Research by the author was supported by the Indian Institute of Technology Bombay post-doctoral research grant.\\ \\
\textbf{Competing Interests} The author has no conflict of interest to declare that is relevant to this article.\\ \\
\textbf{\large Data availability} 
%Data sharing not applicable to this article as no datasets were generated or analysed during the current study.
 No data was used for the research described in the article.
 \fi
\begin{thebibliography}{99}
\bibitem{Ab83}
Abe, E.: Whitehead groups of Chevalley groups over
polynomial rings. Comm. Algebra {\bf 11}(12), 1271--1307 (1983)
\bibitem{ApSt13}
Apte, H., Stepanov, A.: Local-global principle for congruence subgroups of Chevalley groups. Central Eur. J. Math. {\bf 12}, 801--812 (2014)
% %\bibitem{Ba81}
% %A. Bak, {\sl $\K$-theory of forms.}, Ann. of Math. Stud., No. 98. Princeton University Press, Princeton, N. J. University of Tokyo Press, Tokyo, 1981.
% %\bibitem{BBR10}
% %A. Bak, R. Basu and Ravi A. Rao, {\sl Local-global principle for transvection groups}, Proc. Amer. Math. Soc. {\bf 138} (2010), no.4, 1191–1204. 
% %\bibitem{BPT03}
% %A. Bak, V. Petrov and G. Tang, {\sl Stability for Quadratic $\K_{1}$}, $\K$-Theory {\bf 29} (2003), 1--11.

% \bibitem{Ba64}
% H. Bass, {\it K-theory and stable algebra}, Publ. Math. Inst. Hautes $\acute{E}$tudes Sci. {\bf 22} (1964), 5--60.

% \bibitem{Ba73}
% H. Bass, {\sl Unitary algebraic $K$-theory}, Lecture Notes in Mathematics, {\bf 343}, 1973, pp. 57--265.
%\bibitem{BRK05}
%R. Basu, Ravi A. Rao and R. Khanna {\sl On Quillen's local global principle}, Contemporary Mathematics {\bf 390} (2005), Providence, RI; American Mathematical Society. 
\bibitem{BaSi20}
Basu, R., Singh, M.K.: Quillen-Suslin theory for classical groups: revisited over graded rings. Contemp. Math. {\bf 751}, 5--18 (2020)
%\bibitem{Ch82}
%L. G. Chouinard, {\sl Projective modules over Krull semigroups}, Michigan Math. J. {\bf 29}, (1982), 143--148.
\bibitem{GMV91} 
Grunewald, F., Mennicke, J., Vaserstein, L.N.: On symplectic groups over polynomial rings. Math. Zeitschrift {\bf 206}, 35--56 (1991)
% \bibitem{Gu88}
% J. Gubeladze, {\sl Anderson’s conjecture and the maximal class of monoid over which projective modules are free\/}, Math. USSR-Sb. {\bf 63} (1988), 165--188. 
% %\bibitem{Gu90}
% %J. Gubeladze, {\sl Classical algebraic $K$-theory of monoid algebras\/}, Lect. Notes Math., {\bf 437}, (1990), Springer, 36--94. 
\bibitem{Gu92}
Gubeladze, J.: The elementary action on unimodular rows over a monoid ring. J. Algebra {\bf 148}, 135--161 (1992)
% %\bibitem{Gu93}
% %J. Gubeladge, {\it The elementary action on unimodular rows over a monoid ring II}, J. Algebra {\bf 155} (1993), 171--194.
\bibitem{HaMe89} 
Hahn, A.J., O'Meara, O.T.: The classical groups and $\K$-theory. Springer, Berlin (1989)
%\bibitem{Ho95} 
%Y. Hong, {\sl Stabilization of unitary groups over polynomial rings\/}, Chin. Ann. of Math. Ser. B {\bf 16} (1995), no. 2, 177–-190.
%\bibitem{HaVa03}
%R Hazrat and N Vavilov {\sl $K_{1}$ of Chevalley groups are nilpotent\/}, 2003

% \bibitem{Ke87}
% F, Keune {\sl The relativisation of $\K_{2}$\/}, J. Algebra 54 (1) (1987) 159–177.

% \bibitem{Ko78} 
% V.I. Kopeiko, {\sl The stabilization of symplectic groups over a polynomial ring\/},  Math. USSR. Sbornik, {\bf 34} (1978), no. 5, 655–-669.
\bibitem{Lambook}
Lam, T.Y.: Serre's problem on projective modules. Springer Monographs in Mathematics, Springer-Verlag, Berlin (2006)
\bibitem{La18}
Lavrenov, A.: A local–global principle for symplectic $\K_{2}$. Doc. Math. {\bf 23}, 653--675 (2018)
\bibitem{LaSi17} 
Lavrenov, A., Sinchuk, S.: On centrality of even orthogonal $\K_{2}$. J. of Pure and Appl. Algebra {\bf 221}(2), 1134--1145 (2017)

\bibitem{Qu76}
Quillen, D.: Projective modules over polynomial rings. Invent. Math. {\bf 36}, 167--171 (1976)
%\bibitem{St14}
%A. Stavrova, {\sl Homotopy invariance of non-stable K1-functors}{J K-theory} 2014
%\bibitem{St22}
%A. Stavrova, {\sl A1-invariance of non-stable K1-functors in the equicharacteristic case}{Indagationes Mathematicae} 2022

%\bibitem{PeSt08}
%V. PETROV AND A. STAVROVA {\sl ELEMENTARY SUBGROUPS OF ISOTROPIC REDUCTIVE GROUPS} 2008

\bibitem{Si16}
Sinchuk, S.: On centrality of $\K_{2}$ for Chevalley groups of type $\E_{l}$. J. of Pure and Appl. Algebra {\bf 220}(2), 857--875 (2016)
\bibitem{Si25}
Sinchuk, S.: On the $\mathbb{A}^{1}$-invariance of $\K_{2}$ of Chevalley groups of simply-laced type. Eur. J. Math. {\bf 11}(14), 1--42 (2025)
\bibitem{St68}
Steinberg, R.: Lectures on Chevalley groups. Yale University (1968)
% %\bibitem{Su84}
% %A.A. Suslin, {\sl Cancellation over affine varieties} J.Sov. Math. {\bf 27} (1984), 2974--2980.
% %\bibitem{Su77a}
% %A.A. Suslin, {\sl A cancellation theorem for projective modules over affine algebras,} Sov. Math. Dokl. {\bf 18} (1977), 1281--1284.
\bibitem{Su77}
Suslin, A.A.: On the structure of the special linear group over polynomial rings. Math. USSR-Izv. {\bf 11}, 221--238 (1977)
% %\bibitem{Su76}
% %A.A. Suslin, {\sl Projective modules over a polynomial ring are free} Sov. Math. Dokl. {\bf 17} (1976), 1160--1164.
\bibitem{SuKo82}
Suslin, A.A., Kopeiko, V.I.: Quadratic modules and orthogonal group over polynomial rings. J. Sov. Math. {\bf 20}(6), 2665--2691 (1982)
%\bibitem{Ta83}
%G. Taddei, Normalité des groupes élémentaires dans les groupes de Chevalley sur un anneau, in: Applications of Algebraic K-theory to Algebraic Geometry and Number Theory, Part II, Boulder, Colo, 1983, in: Contemp. Math., vol. 55, 1986, pp. 693–710.
\bibitem{Tu83}
Tulenbaev, M.S.: The Steinberg group of a polynomial ring. Math. USSR Sb. {\bf 45}(1), 139--145 (1983)

% % \bibitem{VaSu76}
% % L.N. Vaserstein and A.A. Suslin, {\sl Serre’s problem on projective modules over polynomial rings and algebraic K-theory,} Math. USSR, Izvestija. {\bf 10} (1976), no 5, 937--1001.

% %\bibitem{Va69}
% %L.N. Vaserstein, {\it On the stabilization of the general linear group over a ring}, Math. USSR, Sb. {\bf 8} (1969), no. 3, 383--400.
% \bibitem{Va70}
% L.N. Vaserstein, {\it Stabilization of unitary and orthogonal groups over a ring with involution}, Math. USSR, Sb. {\bf 10} (1970), no. 3, 307--326.
% %\bibitem{Va71}
% %L.N. Vaserstein, {\it Stable range of rings and dimensionality of topological spaces}, Funct. Anal. Appl {\bf 2} (1971), no. 5, 102--110.
\bibitem{VaPl96}
Vavilov, N.A., Plotkin, E.: Chevalley groups over commutative rings: I. Elementary calculations. Acta Appl. Math. {\bf 45}(1), 73--113 (1996)
\bibitem{We10}
Wendt, W.: $\mathbb{A}^{1}$-homotopy of Chevalley groups. J. $\K$-Theory {\bf 5}, 245--287 (2010)
\end{thebibliography}

\end{document}

